\section{Guarantees and Feasibility}
\label{sec:theory}

This section separates what can be guaranteed from what must be measured. All
statements concern the bounded pairwise utility in \cref{sec:problem}; none is
an AUC confidence interval.

\subsection{Why training-only selection is insufficient}

\begin{theorem}[Training-only non-identifiability]
\label{thm:nonidentifiability}
Fix any training transcript $R$ and public inputs. Over an unrestricted family
of target-edge distributions, no selector that observes only $(R,X)$ can be
both nonvacuous and guarantee that the selected structural branch has
nonnegative target utility relative to the public branch.
\end{theorem}

\emph{Proof sketch.}
Construct two target worlds with identical training data, public inputs, and
therefore identical $R$, but reverse the probability mass of candidate pairs
on which $s_R$ and $s_0$ disagree. Any training-only selector has the same
decision distribution in both worlds, while the utility difference changes
sign. A distribution-free no-harm rule must therefore abstain whenever the
two scorers can disagree. The complete construction is in
\cref{app:nonidentifiability}.

The theorem does not say that metadata-based selectors are useless on a
restricted graph family. It says that their safety is an empirical
generalization assumption. \method{} instead obtains direct, private evidence
from the target domain.

\subsection{Transcript privacy}

\begin{theorem}[Adaptive role-wise transcript privacy]
\label{thm:privacy}
Suppose the public role assignment is edge-stable, $\cM_T$ is
$\rho_T(\lambda)$-RDP on $E^T$, and, uniformly for every fixed training
release $r$, $\cM_C(\cdot;r)$ is $\rho_C(\lambda)$-RDP on $E^C$. If inference
uses only released state and public inputs, the joint transcript is
\begin{equation}
 \rho_{\mathrm{joint}}(\lambda)
 =\max\{\rho_T(\lambda),\rho_C(\lambda)\}\text{-RDP}.
\end{equation}
Without an edge-stable disjointness proof, the conservative bound
$\rho_{\mathrm{joint}}\leq\rho_T+\rho_C$ is valid.
\end{theorem}

The proof conditions on the one role containing the neighboring edge.
Post-processing handles certification after a training-edge change, while
uniform conditional DP handles a certification-edge change. R5 uses the
sequential fallback because its historical benchmark split was not designed
to prove edge stability.

\subsection{Finite-population no-harm certificate}

Let $H$ be the sealed positive holdout of size $N$ and
\begin{equation}
 \Delta_H(R)=\frac{1}{N}\sum_{e\in H}d_R(e).
\end{equation}
An edge-keyed public random hash assigns each $e\in H$ independently to
certification with probability $p_C$; $Q5$ is the complement.

\begin{theorem}[Finite-holdout no-harm]
\label{thm:noharm}
Conditional on $(R,H,n_C)$, the certification subset is a simple random sample
without replacement. With the certificate in \cref{eq:certificate},
\begin{equation}
\Pr\!\left(A=1\ \wedge\ \Delta_H(R)<\gamma\mid R,H\right)
\leq \beta_S+\beta_n+\beta_{\mathrm{stat}}.
\end{equation}
\end{theorem}

The edge-keyed hash is essential: adding an edge does not reassign existing
edges. Conditional on the realized bounded values, shared graph endpoints do
not alter the random-sampling argument. A one-sided Serfling bound
\cite{serfling1974probability} gives the
sampling penalty in \cref{eq:certificate}; R5 conservatively sets its
finite-population correction to one. The result certifies the registered
finite holdout, not future edges or another graph.

\subsection{When private certification is possible}

Write $\alpha=\Delta-\gamma>0$ for the effect above the material threshold,
$\chi\geq1$ for an effective dependence factor, and $\nu$ for aggregate
Gaussian noise. The sufficient boundary derived in
\cref{app:feasibility} has the principal form
\begin{equation}
 n=O\!\left(
 \frac{\chi\log(1/\beta)}{\alpha^2}
 {}+\frac{\nu\sqrt{\log(1/\beta)}}{\alpha}
 {}+n_{\min}
 \right).
\label{eq:sufficient-rate}
\end{equation}
Since Gaussian calibration gives
$\nu=\widetilde O(1/\eps)$ under ideal aggregation, the privacy term is
$\widetilde O(1/(\eps\alpha))$. Visible independent client messages replace
$\nu$ by $\sqrt{K}\nu$ in this implementation; this is a utility cost, not
extra privacy composition.

\begin{theorem}[Necessary count; informal]
\label{thm:lower}
For block-replicated Bernoulli hard instances with dependence $\chi$, every
central $(\eps,\delta)$-DP selector with fixed testing error requires, up to
constants and approximate-DP logarithms,
\begin{equation}
 n=\Omega\!\left(\frac{\chi}{\alpha^2}
                 {}+\frac{1}{\eps\alpha}\right).
\end{equation}
\end{theorem}

Thus the sufficient and necessary boundaries match in their principal
dependence on evidence, privacy, and effect size. This does not make the
specific Gaussian certificate minimax optimal in $\delta$, nor does it prove
that arbitrary graph edges follow the hard-instance model. The lower bound
uses the private Le Cam inequality of \cite{acharya2021private}.
